\subsection{交换图}

设，例如，B、C、D、E、F 为五个集合，并设 \(f\) 是从 E 到 F 的映射， \(g\) 是从 B 到 C 的映射， \(h\) 是从 D 到 E 的映射， \(u\) 是从 B 到 D 的映射，且 \(v\) 是从C到E的一个映射。为了概括这类情形，人们常借助图表；例如，前述情形可用下图（E，II，第14页）来概括：

\[
\begin{array}{c} \mathrm{B} \xrightarrow {g} \mathrm{C} \\ u \Big \downarrow \\ \mathrm{D} \xrightarrow [ h ]{} \mathrm{E} \xrightarrow [ f ]{} \mathrm{F}. \end{array}\tag{1}
\]

在这样的图中，符号组 E \(\xrightarrow{f}\) F 示意了 f 是从 E 到 F 的映射这一事实。当对 f 不存在歧义时，可省略字母 f，直接写作 E \(\rightarrow\) F.

当B、C、D、E、F是群（相应地，A-模）且f、g、h、u、v是群（相应地，A-模）的同态时，为简洁起见，称图(1)为群（相应地，A-模）的图。

原则上，图不是数学对象，而仅仅是一个图形，旨在便于阅读证明。在实践中，人们常将图用作缩写符号，以避免逐一命名所有要考虑的集合和映射；因此，人们说“考虑图(1)”，而不是说：“设B、C、D、E、F为五个集合\ldots{} 且v为从C到E的映射”；例如参见第2节命题1的陈述。

例如，考虑下面的图表：

\[
\begin{array}{c} \mathrm{B} \xrightarrow {f} \mathrm{C} \xrightarrow {g} \mathrm{D} \xrightarrow {h} \mathrm{E} \\ b \Big \downarrow \qquad \qquad c \Big \downarrow \qquad \qquad d \Big \downarrow \qquad e \Big \downarrow \\ \mathrm{B} ^ {\prime} \xrightarrow {f ^ {\prime}} \mathrm{C} ^ {\prime} \xrightarrow {g ^ {\prime}} \mathrm{D} ^ {\prime} \xrightarrow {h ^ {\prime}} \mathrm{E} ^ {\prime}. \end{array}\tag{2}
\]

对于由图中若干段按箭头所示方向经过的线段组成的每条路径，都对应一个从第一段起点所表示集合到末段终点所表示集合的映射，即所经过各段对应映射的复合。对于图中的每个顶点，例如C，约定存在一条仅由C构成的路径，并约定与之对应恒等映射 1 \(_{C}\) .

在（2）中，例如，有三条从B出发、终止于D'的路径；相应的映射为 \(d \circ g \circ f\) , \(g' \circ c \circ f\) 和 \(g' \circ f' \circ b\) 。我们说一个图表是交换的，如果对于图表中具有相同起点和相同终点的每一对路径，对应的两个映射相等；特别地，如果一条路径的终点与其起点重合，则对应的映射必须是恒等映射。

为使图（2）可交换，其充要条件是以下关系式

\[
f ^ {\prime} \circ b = c \circ f, \qquad g ^ {\prime} \circ c = d \circ g, \qquad h ^ {\prime} \circ d = e \circ h;\tag{3}
\]

保持；换言之，从（2）中提取的三个方形图可交换是必要且充分的。事实上，关系式（3）蕴含 \(d \circ g \circ f = g' \circ c \circ f\) 因为 \(d \circ g = g' \circ c\) ，以及 \(g' \circ c \circ f = g' \circ f' \circ b\) ，因为 \(c \circ f = f' \circ b\) 因此，从B出发并以D'结束的三条路径给出相同的映射。类似地可以验证，从B出发并以E'结束的四条路径（分别地，从C出发并以E'结束的三条路径）给出相同的映射。关系式(3)意味着从B（分别地，C、D）出发并以C'（分别地，D'、E'）结束的两条路径给出相同的映射。(2)中所有其他顶点对至多由一条路径连接，因此图(2)确实是可交换的。

在下文中，我们将把对其他类型图表进行类似结果的表述与验证这一任务留给读者。

